\documentclass[11pt,a4paper]{article}
\usepackage[a4paper,margin=2.2cm]{geometry}
\usepackage{fontspec}
\setmainfont{Times New Roman}
\setmonofont{Consolas}[Scale=0.86]
\usepackage{amsmath,amssymb,mathtools}
\usepackage{unicode-math}
\setmathfont{Latin Modern Math}
\usepackage{booktabs}
\usepackage{array}
\usepackage{enumitem}
\usepackage{titlesec}
\usepackage{xcolor}
\usepackage{mdframed}
\usepackage[hidelinks]{hyperref}

\definecolor{acc}{HTML}{7A2E1E}
\definecolor{dim}{HTML}{555555}
\definecolor{boxbg}{HTML}{F4F1EC}

\titleformat{\section}{\large\bfseries\color{acc}}{\thesection.}{0.5em}{}
\titlespacing{\section}{0pt}{1.1em}{0.4em}
\titleformat{\subsection}{\bfseries}{\thesubsection.}{0.5em}{}
\titlespacing{\subsection}{0pt}{0.8em}{0.25em}

\setlist[itemize]{leftmargin=1.2em,itemsep=0.15em,topsep=0.3em}
\setlist[enumerate]{leftmargin=1.4em,itemsep=0.15em,topsep=0.3em}
\renewcommand{\arraystretch}{1.15}
\setlength{\parskip}{0.45em}
\setlength{\parindent}{0pt}

\newcommand{\qq}[1]{\texttt{\small #1}}
\newcommand{\note}[1]{{\color{dim}\small #1}}

% Врезка для вывода или ключевого следствия.
\newmdenv[
  backgroundcolor=boxbg, linewidth=0pt, skipabove=0.7em, skipbelow=0.7em,
  innerleftmargin=0.9em, innerrightmargin=0.9em,
  innertopmargin=0.6em, innerbottommargin=0.6em
]{keybox}

\DeclareMathOperator*{\argmax}{arg\,max}
\DeclareMathOperator*{\argmin}{arg\,min}
\DeclareMathOperator{\KL}{KL}
\newcommand{\R}{\mathbb{R}}
\newcommand{\E}{\mathbb{E}}
\newcommand{\Var}{\operatorname{Var}}

\usepackage{fancyvrb}
\fvset{fontsize=\small,xleftmargin=1em}
\begin{document}

{\Large\bfseries Основа: генерация, префилл, вызов инструментов}

\vspace{-0.5em}
\note{vlm-recipes~· что происходит под остальными тремя гайдами}

\section{Что модель делает на самом деле}

Языковая модель — функция одного вида. На вход последовательность
токенов, на выходе распределение вероятностей над следующим токеном:
\[
p_\theta(\,\cdot \mid x_1,\dots,x_{t-1}) \in \Delta^{|V|},
\]
где $|V|$ — размер словаря, у Qwen порядка 150 тысяч.

Больше она не умеет ничего. Всё остальное — диалог, вызов инструментов,
следование инструкции — надстройка над этим единственным действием.

Генерация есть повторное применение: выбрали токен, дописали к входу,
применили снова. Отсюда главное следствие, которое стоит принять сразу:

\begin{keybox}
Модель не \emph{отвечает} на вопрос. Она \emph{продолжает} текст.
Ощущение ответа возникает потому, что текст оформлен так, что
естественным продолжением оказывается реплика ассистента.
\end{keybox}

\section{Диалог как одна строка}

Ролей «пользователь» и «ассистент» для модели не существует. Существуют
токены. Список реплик превращается в одну строку по правилу, которое
называется chat template и хранится в токенизаторе модели.

У Qwen это выглядит так:

\begin{Verbatim}
<|im_start|>system
Ты помогаешь студенту.<|im_end|>
<|im_start|>user
Сформулируй гипотезу<|im_end|>
<|im_start|>assistant
\end{Verbatim}

Обратите внимание на последнюю строку: она открыта. Реплика ассистента
начата и не закончена, и модели остаётся её продолжить — ровно то, что
она умеет. Служебные токены \qq{<|im\_start|>} и \qq{<|im\_end|>} —
обычные элементы словаря, просто зарезервированные под разметку.

Дописывает эту открывающую строку флаг \qq{add\_generation\_prompt=True}.
Без него текст заканчивается на \qq{<|im\_end|>}, и модель с равным
успехом может начать новую реплику пользователя.

Точный формат различается между семействами, поэтому смотреть надо свой:

\begin{Verbatim}
print(processor.apply_chat_template(messages, tokenize=False,
                                    add_generation_prompt=True))
\end{Verbatim}

\section{Префилл}

Приём следует прямо из устройства генерации. Раз модель продолжает
текст, начало ответа можно задать самому — тогда она продолжит уже
написанное.

Технически: добавляем реплику ассистента с нужным началом и просим
шаблон не закрывать её.

\begin{Verbatim}
messages = [
    {"role": "user",      "content": "Сформулируй гипотезу"},
    {"role": "assistant", "content": "Уточняющий вопрос:"},
]
prompt = processor.apply_chat_template(
    messages, tokenize=False, continue_final_message=True
)
\end{Verbatim}

Получается строка, оканчивающаяся на \qq{...assistant\textbackslash n
Уточняющий вопрос:} — и генерация идёт с этой точки.

\begin{keybox}
Флаг \qq{continue\_final\_message=True} здесь обязателен, и одного
\qq{add\_generation\_prompt=False} недостаточно. Шаблон отрисовывает
каждую реплику как \qq{<|im\_start|>роль ... <|im\_end|>}, то есть
закрыл бы и вашу: получилось бы \qq{Уточняющий вопрос:<|im\_end|>} —
реплика завершена, продолжать нечего, префилл не сработал.
\end{keybox}

Проверяется это на своей модели за секунду:

\begin{Verbatim}
for kw in ({"add_generation_prompt": False}, {"continue_final_message": True}):
    print(kw, repr(proc.apply_chat_template(messages, tokenize=False, **kw))[-60:])
\end{Verbatim}

\subsection{Почему это сильнее инструкции}

Инструкция «задай уточняющий вопрос» попадает в условие $x_{<t}$
наравне с текстом пользователя и конкурирует с ним: модель взвешивает
всё содержимое контекста и \emph{решает}, следовать ли ей.

Префилл не участвует в этом соревновании. Он меняет саму точку, из
которой разворачивается авторегрессия:
\[
y \sim p_\theta(\,\cdot \mid x),
\qquad\text{против}\qquad
y \sim p_\theta(\,\cdot \mid x,\; y_1,\dots,y_k),
\]
где $y_1,\dots,y_k$ заданы вами и уже не подлежат пересмотру. Модель
не решает, задавать ли вопрос, — она его уже задала и достраивает
начатое.

\begin{keybox}
Отсюда область применения: префилл идеален там, где форму ответа
выбирает код, а содержание — модель. Состояние документа известно
программе, она и решает, каким должно быть начало реплики.
\end{keybox}

\subsection{Чем это отличается от системного промпта}

Системный промпт работает не потому, что «системный», а потому что при
инструктивном дообучении в данных систематически шло: описание поведения
в системной части, затем соответствующая ему реплика ассистента. Модель
выучила эту корреляцию. Механизм статистический — оттого систему и можно
перебить последующим текстом.

Префилл статистическим не является вовсе. Он не свидетельство, которое
модель взвешивает, а часть самой последовательности:

\begin{center}
\begin{tabular}{@{}ll@{}}
\toprule
системный промпт & свидетельство, которое модель взвешивает \\
префилл & часть самой последовательности \\
\bottomrule
\end{tabular}
\end{center}

Первое можно перевесить другим свидетельством. Второе перевешивать
нечем: оно не в споре, оно уже случилось.

\subsection{Граница действия}

Префилл фиксирует первые $k$ токенов и ровно на этом заканчивается —
дальше модель снова свободна. Начав «Уточняющий вопрос:», она обязана
продолжить с этой точки, но через десяток токенов может задать вопрос
и тут же сама на него ответить.

Гарантируется начало, а не весь ответ. Поэтому в лестнице приёмов
префилл стоит рядом со структурированным декодированием, а не вместо
него: первое задаёт точку старта, второе ограничивает каждый шаг.

\section{Как получается вызов инструмента}

Здесь чаще всего возникает неверная картина, поэтому скажем прямо:

\begin{keybox}
Модель ничего не вызывает. Она порождает строку определённого вида.
Разбирает эту строку, выполняет функцию и возвращает результат
\textbf{ваш код}.
\end{keybox}

\subsection{Что видит модель}

Описания доступных функций подставляются в промпт — обычно в системную
часть, в формате, на котором модель обучалась. Передаются они отдельным
аргументом, а не текстом:

\begin{Verbatim}
tools = [{
    "type": "function",
    "function": {
        "name": "read_document",
        "description": "Текст раздела работы",
        "parameters": {
            "type": "object",
            "properties": {"section_id": {"type": "string"}},
            "required": ["section_id"],
        },
    },
}]
processor.apply_chat_template(messages, tools=tools, tokenize=False,
                              add_generation_prompt=True)
\end{Verbatim}

Шаблон сам развернёт это в текст. Если у модели в \qq{chat\_template}
нет ветки по \qq{tools}, вызов инструментов у неё не обучен — просить
можно, но формат придётся разбирать самому и надёжность будет ниже.

\subsection{Что порождает модель}

Обычный текст, в котором есть размеченный фрагмент. Единого формата
нет — он задаётся шаблоном конкретной модели. Два распространённых:

\begin{Verbatim}
<tool_call>                              <tool_call>
{"name": "read_document",                <function=read_document>
 "arguments": {"section_id": "3.2"}}     <parameter=section_id>
</tool_call>                             3.2
                                         </parameter>
                                         </function>
                                         </tool_call>
\end{Verbatim}

Слева JSON, справа вложенные XML-теги; Qwen3.5 использует правый.
Какой ждёт ваша модель, написано прямо в системной части промпта после
подстановки \qq{tools} — и это же означает, что обучающие данные должны
быть в том же формате, иначе вы учите модель одному, а разбираете другое.

Никакого особого механизма за этим нет. Модель обучена: увидев описания
функций и запрос, для которого нужен вызов, продолжать текст такой
конструкцией. Это тот же авторегрессионный процесс, что порождает любое
другое предложение.

\subsection{Полный цикл}

\begin{Verbatim}
messages = [{"role": "user", "content": "Прочитай раздел 3.2"}]

while True:
    text = generate(messages, tools=tools)      # шаг модели
    calls = parse_tool_calls(text)              # разбор строки
    if not calls:
        break                                   # модель ответила словами
    messages.append({"role": "assistant", "content": text})
    for call in calls:
        result = TOOLS[call["name"]](**call["arguments"])   # ваш код
        messages.append({"role": "tool", "content": result})
\end{Verbatim}

Результат возвращается в диалог как реплика роли \qq{tool}, и следующий
шаг модели видит его наравне с остальным контекстом. Цикл повторяется,
пока модель не перестанет порождать вызовы.

Отсюда два практических следствия. Первое: остановка — это тоже
поведение модели, и ему нужно учить отдельно, иначе цикл не завершится.
Второе: результат инструмента приходит извне, и обучаться предсказывать
его нельзя — модель начнёт сочинять содержимое документов вместо того,
чтобы их запрашивать. В коллаторе эти позиции маскируются.

\section{Префилл и вызов инструмента вместе}

Приёмы складываются. Если нужно гарантировать, что модель вызовет
инструмент, начало вызова задаётся префиллом:

\begin{Verbatim}
{"role": "assistant", "content": '<tool_call>\n{"name": "'}
\end{Verbatim}

Дальше модели остаётся продолжить именем функции — вызывать или
не вызывать она уже не решает. Дешёвая замена структурированному
декодированию, когда его нет под рукой.

Тот же приём наоборот, для форсирования обычного ответа: префилл
непустым текстом закрывает путь к \qq{<tool\_call>} на первом же шаге.

\section{Рассуждающие модели и блок \qq{<think>}}

У рассуждающих моделей шаблон открывает блок размышления сразу после
служебных токенов ассистента — промпт кончается так:

\begin{Verbatim}
<|im_start|>assistant
<think>
\end{Verbatim}

Модель сначала рассуждает внутри блока, затем закрывает его тегом
\qq{</think>} и только потом пишет ответ. Из этого следуют четыре вещи.

\textbf{Префилл ответа должен сначала закрыть блок.} Дописав
«Уточняющий вопрос:» прямо после \qq{<think>}, вы префиллите
\emph{рассуждение}, а не ответ. Чтобы префиллить ответ:

\begin{Verbatim}
prompt += "\n</think>\n\nУточняющий вопрос:"
\end{Verbatim}

Пустой блок \qq{<think>\textbackslash n\textbackslash n</think>} — это
не мусор, а сигнал «рассуждать не нужно». Шаблон сам подставляет его
в завершённые реплики.

\textbf{В истории диалога рассуждение выбрасывается.} Оно нужно для
текущего шага и на следующем только съест контекст и собьёт модель.
Шаблоны Qwen вырезают его сами из всех реплик, кроме последней.
Рассуждение живёт один ход.

\textbf{В обучающих данных без рассуждений шаблон вставит пустой блок.}
Если этот блок попадёт в градиент, модель усвоит «размышлять не нужно»,
и дообучение под правила поведения заодно отучит её думать. Поэтому
коллатор маскирует \qq{<think>...</think>}: градиент не идёт ни за
рассуждение, ни против него.

\textbf{Лимит токенов уходит в рассуждение.} При \qq{max\_new\_tokens=160}
модель успевает только начать думать --- по-английски, независимо
от языка вопроса, --- и ответ не начинается вовсе. Это не «модель
ответила на английском», а обрезанное рассуждение. Штатный выключатель ---
аргумент шаблона:

\begin{Verbatim}
processor.apply_chat_template(messages, add_generation_prompt=True,
                              enable_thinking=False)
\end{Verbatim}

Шаблон дописывает пустой блок сам, и промпт кончается ровно тем, что
стоит перед ответом в обучающих данных: инференс и обучение видят один
и тот же префикс. В ноутбуках и замерах рассуждение выключено именно так.
Второй способ --- дописать \qq{</think>} к промпту после шаблона: это
префилл ответа, а не аргумент шаблона, и он нужен, когда рассуждение
включено, но ответ должен начаться с заданной формы.

\section{Выбор токена}

Распределение получено — остаётся выбрать из него токен.

\begin{center}
\begin{tabular}{@{}lp{9.6cm}@{}}
\toprule
\textbf{Способ} & \textbf{Когда} \\
\midrule
жадный, \qq{temperature=0} & разбор структуры, вызовы инструментов,
распознавание текста — везде, где ответ должен быть воспроизводим \\
\qq{temperature} 0.6--0.8 & свободный диалог; ниже ответы однообразны,
выше теряется связность \\
\qq{top\_p} 0.9 & отсекает хвост распределения, обычно вместе
с температурой \\
\bottomrule
\end{tabular}
\end{center}

Для агентского цикла температура нулевая по причине более важной, чем
воспроизводимость: разбор вызова ломается от одного лишнего символа,
и случайность здесь работает исключительно против вас.

\section{Куда это ведёт}

Всё, что разбирается в остальных трёх гайдах, стоит на этой основе.

Дообучение (\hbox{02-sft-math}) меняет распределение
$p_\theta(\cdot\mid x)$ — какие продолжения модель считает вероятными.
Маскирование определяет, за какие позиции она получает градиент.

Выравнивание (\hbox{03-alignment}) сравнивает два продолжения одного
и того же $x$ и сдвигает распределение в пользу одного из них.

Векторы управления вмешиваются между входом и распределением, сдвигая
активации на промежуточном слое.

Промпт, префилл и грамматика при декодировании не трогают $\theta$
вовсе: первый меняет условие, второй — точку старта, третья — множество
допустимых токенов на каждом шаге. Поэтому они бесплатны и поэтому
пробовать их следует раньше обучения.

\end{document}
